\documentclass[11pt,a4paper]{article}
\usepackage[utf8]{inputenc}
\usepackage[T1]{fontenc}
\usepackage[brazil]{babel}
\usepackage{lmodern}
\usepackage[margin=2.2cm]{geometry}
\usepackage{array,booktabs,longtable}
\usepackage{listings}
\usepackage{xcolor}
\usepackage{xurl}
\usepackage[hidelinks]{hyperref}
\urlstyle{tt}
\setlength{\parindent}{0pt}
\setlength{\parskip}{0.5em}
\setlength{\emergencystretch}{2em}
\setlength{\tabcolsep}{4pt}
\renewcommand{\arraystretch}{1.2}
\lstset{basicstyle=\ttfamily\footnotesize,breaklines=true,breakatwhitespace=false,columns=fullflexible,keepspaces=true,showstringspaces=false,frame=single,rulecolor=\color{black!25},backgroundcolor=\color{black!3},aboveskip=0.7em,belowskip=0.7em}
\hypersetup{pdftitle={Relatório do EP2 -- Classificação de clareza de respostas e-SIC},pdfsubject={ACH2118 -- Introdução ao Processamento de Língua Natural}}
\begin{document}
\begin{center}
{\LARGE\bfseries Relatório do EP2\par}
\vspace{0.3em}
{\large Classificação de clareza de respostas e-SIC\par}
\vspace{0.7em}
\textbf{Disciplina:} ACH2118 -- Introdução ao Processamento de Língua Natural\\
\textbf{Professor:} Ivandré Paraboni\\
\textbf{Projeto:} stack-ivandre
\end{center}
\vspace{0.5em}

\section{Integrantes efetivos do grupo}

\begingroup\small
\begin{longtable}{@{}>{\raggedright\arraybackslash}p{0.57\linewidth}>{\raggedright\arraybackslash}p{0.38\linewidth}@{}}
\toprule
\textbf{Nome completo} & \textbf{Número USP}\\
\midrule
Aaron Ferraz do Amaral Martins da Silva & 14554474\\
Geovanny Luan Piedade Barcelos & 14573909\\
Valéria Caroline Silva & 14677401\\
Kauê Camargo Marinho & 14677422\\
\bottomrule
\end{longtable}
\endgroup

\section{Introdução e estratégia escolhida}

O objetivo é classificar a clareza de respostas a pedidos de informação do e-SIC. A entrada é o texto da resposta (\nolinkurl{resp_text}) e a saída é um dos três rótulos da coluna \nolinkurl{clarity}: \nolinkurl{c1} (nota 1), \nolinkurl{c234} (notas 2 a 4) ou \nolinkurl{c5} (nota 5).

O EP1 combinava BERTimbau e TF-IDF de palavras por média de probabilidades. O EP2 acrescenta um classificador de caracteres e aprende a combinação dos três modelos por \textbf{stacking}. Cada componente estima três probabilidades. A concatenação dessas nove entradas alimenta uma regressão logística multinomial, que atribui a classe de maior probabilidade.

\section{Dados e pré-processamento}

O treino contém 20.092 respostas, com 6.347 exemplos de \nolinkurl{c1}, 6.853 de \nolinkurl{c234} e 6.892 de \nolinkurl{c5}. As divisões agrupam textos iguais após normalização Unicode NFC, \nolinkurl{casefold} e colapso de espaços. Dentro da parcela destinada ao ajuste dos componentes, mantém-se a primeira ocorrência de cada par de texto exato e rótulo. Textos iguais com rótulos diferentes são preservados.

BERTimbau recebe o texto original, com preenchimento dinâmico e limite de 512 tokens, incluindo os especiais. Esse limite afeta 2.174 respostas do treino (10,82\%). Os modelos lexicais recebem o texto inteiro em minúsculas, com frequência sublinear, IDF suavizado e normalização L2. A entrada se restringe à resposta textual. O pedido original e o conteúdo dos anexos ficam fora da representação.

\section{Parâmetros e alternativas avaliados}

Foram comparados os três classificadores individuais, a média 50/50 do EP1, médias duplas com pesos selecionados, médias duplas com calibração, média tripla e stacking com dois ou três componentes. Nas médias duplas, o peso de BERTimbau foi escolhido entre 0, 0,25, 0,50, 0,75 e 1. A calibração ajustou temperaturas entre 0,05 e 20 por perda logarítmica. A configuração de cada componente foi mantida fixa entre as comparações.

\section{Configuração final escolhida}

\begingroup\small
\begin{longtable}{@{}>{\raggedright\arraybackslash}p{0.26\linewidth}>{\raggedright\arraybackslash}p{0.69\linewidth}@{}}
\toprule
\textbf{Componente} & \textbf{Configuração}\\
\midrule
\endfirsthead
\toprule
\textbf{Componente} & \textbf{Configuração}\\
\midrule
\endhead
\bottomrule
\endfoot
BERTimbau & Base cased, seed 20260941, AdamW com taxa de aprendizado $2\times10^{-5}$ e decaimento de pesos 0,01. Lote efetivo 8, 3.557 atualizações, agenda linear com horizonte de 4.702 passos e aquecimento de 470.\\
TF-IDF de palavras & Unigramas e bigramas. Regressão logística L2 com $C=0{,}5$.\\
TF-IDF de caracteres & N-gramas de 3 a 5 caracteres dentro das palavras (\nolinkurl{char_wb}). Regressão logística L2 com $C=1$.\\
Vocabulário lexical & Frequência documental mínima de 3 e até 100.000 características por componente.\\
Stacking triplo & Regressão logística multinomial L2 com $C=1$ e nove probabilidades brutas na ordem BERTimbau, palavras e caracteres.\\
\end{longtable}
\endgroup

\section{Procedimento experimental}

\subsection{Separação dos dados e validação}

O desenvolvimento usou 12 rodadas de validação estratificada por grupos com cinco folds. Em cada fold, cerca de 64\% dos dados foram reservados ao ajuste dos componentes, 16\% à combinação e 20\% à avaliação. Os grupos foram mantidos separados entre as três parcelas. Os pesos e temperaturas das alternativas foram selecionados em três subdivisões da parcela de combinação. Cada linha integrou a avaliação uma vez por rodada. As médias entre rodadas são descritivas, pois as rodadas reutilizam as mesmas respostas.

\subsection{Ajuste final e previsão do teste}

Os três componentes finais foram ajustados em 15.035 respostas após deduplicação. O combinador foi ajustado em outras 4.101 respostas de grupos distintos, usando as probabilidades dos componentes já treinados. O pacote final gerou os 897 rótulos do teste 2, cujo gabarito não está disponível.

\section{Resultados do modelo final escolhido}

\subsection{Desenvolvimento}

\begingroup\small
\begin{longtable}{@{}lrr@{}}
\toprule
\textbf{Modelo} & \textbf{Acurácia média} & \textbf{Macro F1 médio}\\
\midrule
BERTimbau & 45,0876\% & 44,9667\%\\
TF-IDF de palavras & 44,7948\% & 44,4939\%\\
TF-IDF de caracteres & 44,4634\% & 44,2868\%\\
Média 50/50 (estratégia do EP1) & 45,9329\% & 45,7694\%\\
Média tripla & 46,1436\% & 45,9378\%\\
Stacking duplo & 46,1282\% & 45,7854\%\\
\textbf{Stacking triplo} & \textbf{46,1830\%} & 45,8372\%\\
\bottomrule
\end{longtable}
\endgroup

O stacking triplo obteve a maior acurácia média, com vantagem de 0,2501 ponto percentual sobre a média 50/50 e melhora em 11 das 12 rodadas. A média tripla obteve o maior Macro F1. O stacking também apresentou menor erro de calibração que a média 50/50 (ECE de 0,01370 ante 0,04529).

\subsection{Comparação com o EP1 no teste 1}

O pacote final foi comparado com as previsões originalmente entregues no EP1 nas mesmas 900 respostas, com 300 exemplos de cada classe.

\begingroup\small
\begin{longtable}{@{}lrrr@{}}
\toprule
\textbf{Modelo} & \textbf{Acertos} & \textbf{Acurácia} & \textbf{Macro F1}\\
\midrule
EP1 & 443/900 & 49,22\% & 49,15\%\\
EP2 & 462/900 & 51,33\% & 50,81\%\\
\bottomrule
\end{longtable}
\endgroup

Os dois modelos acertaram 378 respostas e erraram 373. Apenas o EP1 acertou 65 respostas e apenas o EP2 acertou 84. A melhora foi de 2,11 pontos percentuais. O teste de McNemar exato bilateral resultou em $p=0{,}1401$ (estatística descritiva sem correção de 2,4228). A diferença observada ficou aquém do nível de significância de 5\%.

O recall de \nolinkurl{c1} passou de 47,33\% para 55,67\%, e o de \nolinkurl{c5} de 56,67\% para 61,33\%. Para \nolinkurl{c234}, caiu de 43,67\% para 37,00\%. O teste já havia sido usado em avaliações anteriores. Há 88 respostas idênticas a textos de treino, 105 após normalização, além de duplicatas e rótulos conflitantes no próprio teste. Esses fatores limitam a interpretação da comparação como evidência de generalização.

\subsection{Interpretação e análise de erros}

No modelo de palavras da primeira divisão, os termos \emph{demanda} e \emph{não se} têm coeficientes positivos para \nolinkurl{c1}, \emph{sistemas} e \emph{concurso} para \nolinkurl{c234}, e \emph{dnpm} e \emph{questionário} para \nolinkurl{c5}. Esses pesos refletem associações no conjunto de treino, inclusive vocabulário específico de órgãos e procedimentos.

A contribuição local de cada termo resulta do produto entre o TF-IDF e o coeficiente da classe. Em uma resposta que remetia a anexos (linha 6636, com índice iniciado em zero), \emph{encaminhamos} contribuiu com $+0{,}1179$ para o escore de \nolinkurl{c5}. O classificador de palavras acertou a classe, enquanto a combinação selecionada naquele fold previu \nolinkurl{c234}. Em outra resposta, uma justificativa de sigilo acompanhava informação parcial. A combinação previu \nolinkurl{c1} para um exemplo de \nolinkurl{c234}. O texto excedia 512 tokens, e o pedido original seria necessário para avaliar a adequação da resposta.

\section{Repositório de código}

O código, os modelos ajustados, a configuração, as planilhas e as métricas agregadas estão disponíveis em \url{https://github.com/Zennitte/stack-ivandre}. As instruções de instalação, inferência e treinamento estão no \nolinkurl{README.md}.

\section{Instruções de reprodução}

Com Python 3.12 e Git LFS, obtenha os pesos e instale as dependências. O exemplo usa o índice de rodas PyTorch para CPU.

\begin{lstlisting}
git lfs pull
python -m pip install -r requirements.txt --extra-index-url https://download.pytorch.org/whl/cpu
python predict.py --input data/test.xlsx --sheet test2 --output predictions/nova_execucao.xlsx --device cpu
\end{lstlisting}

A saída preserva as linhas e o conteúdo original da planilha e preenche \nolinkurl{clarity}. O código de ajuste final, o split e a revisão do pré-treinado constam no repositório.

\end{document}
